\documentclass[11pt,a4paper]{article}
\usepackage[utf8]{inputenc}
\usepackage{amsmath,amssymb,amsthm}
\usepackage{mathrsfs}
\usepackage{geometry}
\geometry{margin=1in}
\usepackage{hyperref}
\usepackage{enumitem}

\newtheorem{definition}{Definition}[section]
\newtheorem{theorem}{Theorem}[section]
\newtheorem{proposition}{Proposition}[section]
\newtheorem{remark}{Remark}[section]

\title{\textbf{Classical Field Theory\\in the Algebraic Framework}\\[6pt]
\large Lecture by Romeo Brunetti}
\author{Lecture Notes from the AQFT 50th Anniversary Conference}
\date{}

\begin{document}
\maketitle

\begin{abstract}
Romeo Brunetti presents a rigorous algebraic formulation of classical interacting
field theory on globally hyperbolic spacetimes, developed in collaboration with
Fredenhagen and Rejzner. The framework distinguishes kinematical structure
(configuration space, observables, wavefront set conditions) from dynamical
structure (generalized Lagrangians, M\o ller operators, Peierls brackets).
The resulting observable algebra is nuclear, carries a Poisson structure via
Peierls brackets, and fits naturally into the locally covariant framework.
\end{abstract}

\tableofcontents

%---------------------------------------------------------------------
\section{Motivation}
%---------------------------------------------------------------------

Classical field theory in the algebraic setting received limited attention
historically (Roberts--Leyland 1978, D\"utsch--Fredenhagen 2003). Physicists
typically use either geometric methods (multisymplectic geometry, jet bundles)
or functional methods (infinite-dimensional Lagrangian mechanics). The algebraic
perspective offers a fresh viewpoint---one that is strongly biased by structures
appearing in perturbative quantum field theory.

A key insight: \emph{classical field theory, when formulated algebraically,
is not so drastically infinite-dimensional---the observable algebra is nuclear,
meaning it is very close to being finite-dimensional.}

%---------------------------------------------------------------------
\section{Kinematical Structure}
%---------------------------------------------------------------------

\subsection{Configuration Space}

The configuration space is the space of all smooth functions on a globally
hyperbolic spacetime $(M, g)$:
\begin{equation}
  \mathcal{E} = C^\infty(M),
\end{equation}
equipped with the standard Fr\'echet topology.

\subsection{Observables as Nonlinear Functionals}

Observables are smooth functionals $F: \mathcal{E} \to \mathbb{R}$
subject to two types of restrictions:

\subsubsection{Support Properties}

The support of a functional $F$ is defined analogously to distribution theory:
\begin{equation}
  \mathrm{supp}(F) = M \setminus \{x \in M \mid F \text{ is constant
  on configurations differing only near } x\}.
\end{equation}
We require $\mathrm{supp}(F)$ to be compact.

\subsubsection{Additivity}

A functional $F$ is \textbf{additive} if
\begin{equation}
  F(\varphi_1 + \varphi_2 + \varphi_3) = F(\varphi_1 + \varphi_2)
  + F(\varphi_2 + \varphi_3) - F(\varphi_2)
\end{equation}
whenever $\mathrm{supp}(\varphi_1) \cap \mathrm{supp}(\varphi_3) = \emptyset$.
This generalizes the shift-like properties of distributions to nonlinear
functionals, allowing decomposition into arbitrarily small pieces.

\begin{proposition}
Any additive, compactly supported functional can be decomposed into a
\emph{finite sum} of such functionals with arbitrarily small support.
\end{proposition}

\subsection{Regularity: Differentiability in Fr\'echet Spaces}

The directional derivative of $F$ at $\varphi$ in direction $h$ is:
\begin{equation}
  F^{(1)}(\varphi)[h] = \lim_{t \to 0} \frac{F(\varphi + th) - F(\varphi)}{t}.
\end{equation}
$F$ is $C^1$ if this derivative exists for all $\varphi, h$ and the map
$(\varphi, h) \mapsto F^{(1)}(\varphi)[h]$ is jointly continuous. Higher
derivatives are defined iteratively. The fundamental results of basic
calculus (Leibniz rule, chain rule, fundamental theorem, Schwarz lemma)
remain valid.

\subsection{Local and Microlocal Functionals}

\begin{definition}[Local Functional]
A smooth functional $F$ is \textbf{local} if the support of any
$n$-th derivative $F^{(n)}(\varphi)$ lies on the thin diagonal
$d_n \subset M^n$, with wavefront set orthogonal to the tangent
bundle of $d_n$.
\end{definition}

Typical example: $F(\varphi) = \int f(x)\, L(\varphi(x), \partial\varphi(x),\ldots)\, dv_g(x)$,
where $L$ is a polynomial in the field and finitely many derivatives.

\begin{definition}[Microlocal Functional]
A smooth functional $F$ is \textbf{microlocal} if the wavefront set
of $F^{(n)}(\varphi)$ does not contain elements where all covectors
simultaneously lie in $\overline{V}_+^*$ or all in $\overline{V}_-^*$.
\end{definition}

\begin{theorem}
In the smooth case, local functionals and additive functionals coincide.
The algebra $\mathscr{F}(M)$ of microlocal functionals is a
\textbf{nuclear, sequentially complete topological algebra}.
\end{theorem}

%---------------------------------------------------------------------
\section{Dynamical Structure}
%---------------------------------------------------------------------

\subsection{Generalized Lagrangians}

\begin{definition}
A \textbf{generalized Lagrangian} is a map $\mathfrak{L}: \mathcal{D}(M) \to
\mathcal{F}_{\mathrm{loc}}(\mathcal{E})$ satisfying:
\begin{enumerate}[nosep]
  \item $\mathrm{supp}(\mathfrak{L}(f)) \subset \mathrm{supp}(f)$.
  \item $\mathfrak{L}(f_1 + f_2) = \mathfrak{L}(f_1) + \mathfrak{L}(f_2)$
        whenever $\mathrm{supp}(f_1) \cap \mathrm{supp}(f_2) = \emptyset$.
\end{enumerate}
\end{definition}

For $f \equiv 1$ on a relatively compact region $\mathcal{O}$, the
Euler--Lagrange equations are:
\begin{equation}
  \frac{\delta \mathfrak{L}(f)}{\delta \varphi}\bigg|_\mathcal{O} = 0.
\end{equation}

\subsection{Linearization and Hyperbolicity}

The linearized field equation around any background $\varphi_0$ is obtained
from the second functional derivative of $\mathfrak{L}$:
\begin{equation}
  P_{\varphi_0} = \frac{\delta^2 \mathfrak{L}(f)}{\delta\varphi^2}\bigg|_{\varphi_0}.
\end{equation}
This is required to be a \textbf{strictly hyperbolic differential operator},
guaranteeing the existence of unique advanced and retarded Green functions
on any globally hyperbolic spacetime.

\subsection{M\o ller Operators}

\begin{definition}
Given two Lagrangians $\mathfrak{L}_1$ and $\mathfrak{L}_2$, a \textbf{M\o ller
operator} is an endomorphism $r: \mathcal{E} \to \mathcal{E}$ satisfying:
\begin{enumerate}[nosep]
  \item \textbf{Intertwining}: $\frac{\delta \mathfrak{L}_2}{\delta\varphi}
        \circ r = \frac{\delta \mathfrak{L}_1}{\delta\varphi}$.
  \item \textbf{Retardation}: $r(\varphi) = \varphi$ outside the future
        of the perturbation support.
\end{enumerate}
\end{definition}

\begin{theorem}[Existence and Uniqueness]
For $\mathfrak{L}_2 = \mathfrak{L}_1 + \lambda V$ with $V$ compactly supported,
the M\o ller operator exists and is unique in a neighborhood of $\lambda = 0$,
for any background configuration $\varphi_0$.
\end{theorem}

The proof uses the Nash--Moser implicit function theorem, adapted to
Fr\'echet spaces via:
\begin{enumerate}[nosep]
  \item Decomposition of the perturbation into small pieces.
  \item Solution via a priori energy estimates for the propagator.
  \item Composition to glue pieces back together.
\end{enumerate}

\subsection{Peierls Brackets}

The Poisson structure is introduced via the \textbf{Peierls bracket}:
\begin{equation}
  \{F, G\}_{\mathfrak{L}} = \left\langle \frac{\delta F}{\delta\varphi},\;
  (G^+ - G^-) \frac{\delta G}{\delta\varphi} \right\rangle,
\end{equation}
where $G^\pm$ are the advanced and retarded Green functions of the
linearized equation around $\varphi$.

\begin{theorem}
The Peierls bracket satisfies the Jacobi identity and the Leibniz rule,
making $(\mathscr{F}(M), \mathfrak{L}, \{\cdot,\cdot\}_{\mathfrak{L}})$
a \textbf{Poisson algebra}.
\end{theorem}

%---------------------------------------------------------------------
\section{Consequences}
%---------------------------------------------------------------------

\subsection{Darboux Theorem}

\begin{theorem}[Functional Darboux Theorem]
The M\o ller operators are \textbf{canonical transformations}: they intertwine
the Peierls brackets of $\mathfrak{L}_1$ and $\mathfrak{L}_2$. In particular,
the Poisson bracket can be locally put into normal form.
\end{theorem}

\subsection{On-Shell Theory}

The space of functionals vanishing on solutions forms a \textbf{Poisson ideal}.
The quotient $\mathscr{F}(M)/\mathcal{J}$ gives the on-shell observable
algebra. In the nonlinear case, solutions may have finite lifespan, so
the morphisms between quotient algebras may \emph{fail to be injective}---a
genuinely new phenomenon.

\subsection{Symplectic Structure}

The Casimir functionals (center of the Poisson algebra) are generated by
elements $\int f(\varphi)\,\frac{\delta\mathfrak{L}}{\delta\varphi}\, dv_g$
and constants. Quotienting by the Casimir ideal yields a \textbf{symplectic}
(nondegenerate) Poisson algebra.

All quotients remain in the nuclear topological class.

%---------------------------------------------------------------------
\section{Locally Covariant Formulation}
%---------------------------------------------------------------------

The entire construction is functorial. A natural Lagrangian is a natural
transformation $\mathfrak{L}: \mathcal{D} \Rightarrow \mathscr{F}_{\mathrm{loc}}$
satisfying additivity and naturality:
\begin{equation}
  \alpha_\chi \circ \mathfrak{L}_M = \mathfrak{L}_N \circ \chi_*.
\end{equation}

\begin{theorem}
Additivity in test functions combined with naturality automatically implies
that $\mathfrak{L}_M(f)$ is a local functional. The functor
$\mathscr{F}_{\mathfrak{L}}: \mathbf{Loc} \to \mathbf{NucPois}$ satisfies
all axioms of a locally covariant field theory, including the causality axiom
(now in terms of the Peierls bracket).
\end{theorem}

%---------------------------------------------------------------------
\section{Conclusions and Outlook}
%---------------------------------------------------------------------

\begin{itemize}[nosep]
  \item The algebraic framework provides a rigorous foundation for
        interacting classical field theory with good topological properties.
  \item The construction works for semilinear equations; for quasilinear
        equations, tools of Klainerman may be applicable.
  \item Extension to \textbf{gauge theories} requires BRS/BV constructions.
  \item Extension to \textbf{Einstein's equations} requires paradifferential
        calculus.
  \item A long-term goal: deformation quantization of this classical
        Poisson algebra to produce interacting quantum field theories.
\end{itemize}

\end{document}
